\chapter{Aula 1: TDD sob carga (28/09/2026)}
\label{cap:aula1}

Esta aula transformou um pedido de produto em testes, e os testes em
código, em passos pequenos. O domínio (progressão de cargas no
powerlifting) foi só um pretexto: tinha números, limites e regras que
brigam entre si, e é isso que torna o TDD interessante.

Este capítulo é um resumo para revisar o que foi visto e achar onde
aprofundar. O fio condutor é: pedido de produto, perguntas, exemplos,
testes, código.

\section{Os ciclos do TDD}

Segundo Uncle Bob (\citeauthor{artigo:cycles-tdd}, \citeyear{artigo:cycles-tdd},
\citetitle{artigo:cycles-tdd}), o TDD acontece em quatro ciclos aninhados,
cada um com seu próprio relógio.

\begin{center}
\small
\begin{tabular}{@{}l l p{8cm}@{}}
\toprule
Ciclo & Relógio & O que fazer \\
\midrule
Nano & segundos & As três leis: só escrever código de produção com um teste
falhando, só escrever teste suficiente para falhar, só escrever código
suficiente para passar \\
Micro & minutos & Vermelho, verde, refatorar \\
Milli & $\sim$10 min & Testes cada vez mais específicos, código cada vez
mais genérico \\
Primário & $\sim$1 hora & Parar e olhar as fronteiras da arquitetura \\
\bottomrule
\end{tabular}
\end{center}

Os dois primeiros ciclos são os que você sente no dedo: as três leis
mandam a cadência, e o ritmo vermelho, verde, refatorar fecha cada volta.
O ciclo milli é o de generalização: a cada teste novo o teste fica mais
específico e o código de produção, mais genérico. O ciclo primário é o
momento de levantar a cabeça e olhar a arquitetura, por exemplo onde
ficam as bordas do sistema.

Uma extensão natural da ideia de ``código mais genérico'' é a
\citetitle{artigo:tpp} (\citeauthor{artigo:tpp}, \citeyear{artigo:tpp}),
indicada abaixo em \enquote{onde aprofundar}.

\section{Fake it, triangulação e implementação óbvia}

Para chegar ao verde há dois caminhos, ambos de Kent Beck
(\citetitle{livro:beck2002}):

\begin{description}
\item[Sem saber a solução.] Fingir (\emph{fake it}): devolver a constante
  que o teste espera, por exemplo \texttt{return 100.0}. Em seguida
  \emph{triangular}: escrever um segundo exemplo que a constante não
  satisfaz, o que força a generalizar.
\item[Sabendo a solução.] Escrever direto a \emph{implementação óbvia}.
\end{description}

Ficou vermelho sem entender por quê? O passo foi grande demais. Volte,
desfaça, e dê um passo menor.

\section{Do pedido ao teste}

O caminho completo é:

\begin{center}
frase do ticket $\rightarrow$ perguntas ao produto $\rightarrow$
exemplos aprovados $\rightarrow$ teste com o nome da regra.
\end{center}

A frase do ticket quase nunca basta. As perguntas que mais rendem:

\begin{enumerate}
\item Exatamente no limite conta?
\item Com que precisão?
\item E se o dado vier errado?
\item E se duas regras brigarem?
\item O que nunca pode acontecer?
\end{enumerate}

As respostas viram exemplos concretos, aprovados por quem pediu, e cada
exemplo vira um teste cujo nome é a regra de negócio. O livro
\citetitle{livro:adzic2011} (\citeauthor{livro:adzic2011}) é a referência
para esse passo.

\section{Vocabulário: Fowler e os dublês}

O vocabulário vem de Martin Fowler
(\citetitle{artigo:testdouble}, \citetitle{artigo:mocks-arent-stubs}) e
de \citeauthor{livro:meszaros2007}.

\begin{itemize}
\item \termo{SUT}{\emph{System Under Test}: o trecho de código que o teste
  está exercitando.} (o objeto sob teste) e o \emph{colaborador} (o que o
  SUT usa).
\item \termo{Fixture}{O estado inicial conhecido de que o teste precisa
  para rodar: objetos criados, dados carregados.}: o estado de partida do
  teste.
\item As quatro fases do teste: \emph{setup}, \emph{exercise},
  \emph{verify}, \emph{teardown}.
\end{itemize}

Os cinco dublês de teste (\emph{test doubles}):

\begin{description}
\item[Dummy] preenche um parâmetro e nunca é de fato usado.
\item[Stub]\termo{Stub}{Dublê que devolve respostas prontas às chamadas
  feitas durante o teste.} devolve respostas prontas.
\item[Fake] tem implementação funcional, mas simplificada (por exemplo um
  repositório em memória).
\item[Spy] registra como foi chamado, para o teste verificar
  \emph{estado} depois.
\item[Mock]\termo{Mock}{Dublê programado com as chamadas que espera
  receber; o teste verifica o \emph{comportamento}.} verifica
  \emph{comportamento}: se as chamadas esperadas aconteceram.
\end{description}

Verificar estado (spy) ou comportamento (mock) é uma escolha de projeto.
Na prática do dojo: objetos reais sempre que possível, e dublês só nas
bordas (banco, rede, relógio).

\section{Robustez}

São cinco camadas de confiança, cada uma cobrindo o que a anterior deixa
passar:

\begin{enumerate}
\item exemplos;
\item bordas;
\item entradas inválidas;
\item propriedades (com \citetitle{doc:jqwik});
\item mutação (com \citetitle{doc:pitest}).
\end{enumerate}

Para lembrar o que testar, use o \textbf{Right-BICEP}, do livro
\citetitle{livro:langr}. Critério de aceitação sem teste é critério não
entregue.

A quinta camada é o assunto da aula 2 (capítulo~\ref{cap:aula2}).

\section{Onde aprofundar}

Se for ler um só livro, comece pelo de Kent Beck (TDD) ou pelo
\emph{Accelerate} (visão de entrega, aula 2). Os títulos marcados com
\pt têm edição em português.

\subsection*{Livros}
\small
\begin{longtable}{@{}p{1.5cm} p{4.4cm} p{2.6cm} p{4cm}@{}}
\toprule
Tema & Livro & Autor(es), ano & Por que ler \\
\midrule \endhead
TDD & \citetitle{livro:beck2002} \pt{} \emph{TDD: Desenvolvimento Guiado por
  Testes} & \citeauthor{livro:beck2002}, \citeyear{livro:beck2002}
  & A origem do fake it, da triangulação e da implementação óbvia \\
TDD & \citetitle{livro:martin2021} \pt{} \emph{Craftsmanship Limpo}
  & \citeauthor{livro:martin2021},
  \citeyear{livro:martin2021} & TDD passo a passo com katas em Java e JUnit \\
TDD & \citetitle{livro:martin2008}, cap.~9 \pt{} \emph{Código Limpo}
  & \citeauthor{livro:martin2008}, \citeyear{livro:martin2008}
  & As três leis e o F.I.R.S.T. \\
TDD & \citetitle{livro:martin2011}, cap.~5 \pt{} \emph{O Codificador Limpo}
  & \citeauthor{livro:martin2011}, \citeyear{livro:martin2011}
  & TDD como disciplina profissional \\
TDD & \citetitle{livro:freeman2009} & \citeauthor{livro:freeman2009},
  \citeyear{livro:freeman2009} & TDD de fora para dentro, com mocks nas
  bordas \\
Testes & \citetitle{livro:aniche2015} \pt & \citeauthor{livro:aniche2015},
  \citeyear{livro:aniche2015} (Casa do Código) & Introdução prática em
  português, com Java \\
Testes & \citetitle{livro:aniche2022} & \citeauthor{livro:aniche2022},
  \citeyear{livro:aniche2022} & Testes por especificação, bordas,
  propriedades e mutação, em Java \\
Testes & \citetitle{livro:langr} & \citeauthor{livro:langr}
  & O Right-BICEP da aula \\
Testes & \citetitle{livro:khorikov2020} & \citeauthor{livro:khorikov2020},
  \citeyear{livro:khorikov2020} & Quando usar mocks, estado vs
  comportamento \\
Testes & \citetitle{livro:meszaros2007} & \citeauthor{livro:meszaros2007},
  \citeyear{livro:meszaros2007} & De onde vem o vocabulário dos dublês de
  teste \\
Testes & \citetitle{livro:feathers2004} \pt{} \emph{Trabalho Eficaz com
  Código Legado} & \citeauthor{livro:feathers2004},
  \citeyear{livro:feathers2004} & Como colocar testes em código que não tem \\
Produto & \citetitle{livro:adzic2011} & \citeauthor{livro:adzic2011},
  \citeyear{livro:adzic2011} & Do pedido de produto a exemplos que viram
  testes \\
\bottomrule
\end{longtable}
\normalsize

\subsection*{Blog posts e artigos}
Todos em inglês e gratuitos. Os dois posts do Uncle Bob são a base das
aulas.

\small
\begin{longtable}{@{}p{1.5cm} p{6.5cm} p{4.5cm}@{}}
\toprule
Tema & Artigo & Autor, ano \\
\midrule \endhead
TDD & \citetitle{artigo:cycles-tdd} & \citeauthor{artigo:cycles-tdd},
  \citeyear{artigo:cycles-tdd} \\
TDD & \citetitle{artigo:tpp} & \citeauthor{artigo:tpp},
  \citeyear{artigo:tpp} \\
Dublês & \citetitle{artigo:mocks-arent-stubs} &
  \citeauthor{artigo:mocks-arent-stubs} \\
Dublês & \citetitle{artigo:testdouble} & \citeauthor{artigo:testdouble} \\
\bottomrule
\end{longtable}
\normalsize

\subsection*{Vídeos}
Palestras e tutoriais, a maioria em inglês; os marcados com \pt são em
português. Se tiver tempo para um só, assista ao do Ian Cooper
(\citetitle{video:cooper-tdd-wrong}).

\small
\begin{longtable}{@{}p{1.5cm} p{7cm} p{4cm}@{}}
\toprule
Tema & Vídeo & Quem \\
\midrule \endhead
TDD & \citetitle{video:cooper-tdd-wrong} (Build Stuff 2013) &
  \citeauthor{video:cooper-tdd-wrong} \\
TDD & \citetitle{video:cooper-tdd-revisited} (NDC Porto 2023) &
  \citeauthor{video:cooper-tdd-revisited} \\
TDD & \citetitle{video:martin-clean-code-4} (aula sobre TDD) &
  \citeauthor{video:martin-clean-code-4} \\
TDD & \citetitle{video:beck-theme-variations} &
  \citeauthor{video:beck-theme-variations} \\
TDD & \citetitle{video:tdd-is-dead-2} & Kent Beck, Martin Fowler e David
  Heinemeier Hansson \\
Testes & \citetitle{video:aniche-escrever-testes} \pt &
  \citeauthor{video:aniche-escrever-testes} \\
\bottomrule
\end{longtable}
\normalsize

\subsection*{Ferramentas e documentação}
Nesta aula: \citetitle{doc:junit5}, \citetitle{doc:assertj},
\citetitle{doc:mockito} (testes unitários), \citetitle{doc:jqwik}
(propriedades) e \citetitle{doc:pitest} (mutação). A lista completa, com
links, está no apêndice de ferramentas.

\section{Desafios para casa}

\desafio{Passe um ticket real do seu time pelo caminho frase $\rightarrow$
perguntas $\rightarrow$ exemplos $\rightarrow$ testes. Anote quais
perguntas o ticket não respondia.}

\desafio{Rode o PIT no projeto do dojo e mate os mutantes sobreviventes.}

\desafio{Refaça os katas do \citetitle{livro:martin2021}.}

